\documentclass[11pt]{article}

\usepackage[top=0.7in, bottom=0.7in, left=0.7in, right=0.7in]{geometry}
\usepackage{enumitem}
\setlist{nosep, leftmargin=1.2em}
\usepackage{titlesec}
\titlespacing{\section}{0pt}{12pt}{6pt}
\titlespacing{\subsection}{0pt}{8pt}{4pt}
\usepackage{fontspec}
\setmainfont{Times New Roman}
\usepackage[colorlinks=true, linkcolor=black, urlcolor=blue]{hyperref}
\pagenumbering{gobble}
\setlength{\parindent}{0pt}
\linespread{1.1}

\begin{document}

\begin{center}
{\Large \textbf{Karim Naguib --- Curriculum Vitae}}\\[2mm]
Economist/Data Scientist | Bayesian Methods \& Causal Inference\\[1mm]
karimn2.0@pm.me \quad karimn.co \quad github.com/karimn
\end{center}

\vspace{4mm}

\section*{Summary}

Economist and applied Bayesian statistician with 10+ years building \textbf{hierarchical models} for \textbf{longitudinal and survival data}. Most recently at AstraZeneca, building \textbf{patient-level digital-twin models} of oncology tumor dynamics with real-world evidence integration; earlier work spans field experiments in East Africa and South Asia, real-estate market design, and health-behavior interventions.

\section*{Skills}

\textbf{Methods:} Bayesian hierarchical modeling, state-space and longitudinal models, multistate survival and competing risks, joint longitudinal--survival (``digital twin'') models, cross-validation and predictive assessment (LFO-CV, PSIS-LOO), MCMC diagnostics and prior sensitivity, causal inference and randomized experiment design, Gaussian processes

\textbf{Languages:} R (tidyverse), Stan, Julia, SQL, C/C++

\textbf{Agentic Workflows for Statistical Modeling:} Built and maintain an agentic-workflow layer on Claude Code; at AstraZeneca it was shared across the modeling team via an internal plugin marketplace. Collaborative mathematical derivation with Claude as derivation partner, Stan diagnostic sub-agents, prior sensitivity review, \texttt{targets} pipeline explorer, MCP-integrated skills for Domino job management, reusable plugin and skill templates codifying team conventions for Bayesian workflow.

\textbf{Infrastructure:} CmdStan / cmdstanr, \texttt{targets} pipelines, Domino, Docker, AWS, Databricks, Snowflake, PostgreSQL, git, SLURM, RStudio / VS Code

\section*{Education}

\textbf{Ph.D.\ in Economics} | \textit{Boston University} \hfill Boston, MA | 2014\\
\textit{Fields: development, health, and experimental economics}

\medskip

\textbf{B.S.\ in Computer Science, Minor: Mathematics} | \textit{The American University in Cairo} \hfill Cairo, Egypt | 1999

\section*{Experience}

\textbf{Independent Researcher / Consultant} | \textit{Self-employed} \hfill Jul 2026 -- present
\begin{itemize}
  \item Independent Bayesian modeling and causal-inference work; continued development of agentic tooling for statistical modeling
\end{itemize}

\medskip

\textbf{Senior Data Scientist} | \textit{AstraZeneca} \hfill Aug 2023 -- Jul 2026
\begin{itemize}
  \item Bayesian digital-twin modeling of oncology tumor dynamics for patient-level forecasting of Phase 3 trial outcomes, integrating emerging Phase 2 evidence with historical trial data and real-world data to support Phase 3 initiation decisions (Ph3ID)
  \item Built hierarchical state-space models of longitudinal tumor size (Stein-Fojo framework) with population/trial/patient hierarchy, non-centered parameterization, and Student-t priors to handle prior-data conflict in sparse cohorts; validated via leave-future-out cross-validation
  \item Developed multistate (illness-death) survival models decomposing progression-free survival, jointly fit with the tumor-dynamics model so the patient-level latent tumor state drives burden-related progression risk mechanistically
  \item Contributed modular Stan model architecture and a \texttt{targets}-based pipeline used across multiple trial analyses; routine work on MCMC diagnostics (R-hat, ESS, E-BFMI, divergences) and prior sensitivity
  \item Built and shared an agentic-workflow layer --- collaborative mathematical derivation, diagnostic sub-agents, and MCP-integrated pipelines --- adopted by the modeling team for deriving, implementing, and debugging hierarchical Bayesian models
\end{itemize}

\medskip

\textbf{Independent Researcher} | \textit{Self-employed} \hfill Nov 2022 -- Aug 2023
\begin{itemize}
  \item Built a \href{https://www.github.com/karimn/FundingPOMDPs.jl}{Julia package} for a partially observable MDP solution using Monte Carlo tree search to optimize intervention implementation and evaluation decisions
  \item Bayesian hierarchical models for future-state prediction inside the POMDP; particle filters for faster online planning
  \item \href{https://github.com/karimn/grokking-drl}{Implemented} various deep RL agents and models using Julia and Flux.jl
\end{itemize}

\medskip

\textbf{Senior Data Scientist} | \textit{Opendoor} \hfill May 2021 -- May 2022
\begin{itemize}
  \item Designed randomized experiments to estimate demand and supply elasticities informing pricing policy across regional real-estate markets
  \item Built hierarchical Bayesian models and Gaussian-process priors over spatial/temporal regional structure to pool information across markets while preserving local heterogeneity
  \item Worked across experiment design, inference, and decision-facing deployment
\end{itemize}

\medskip

\textbf{Independent Researcher / Consultant} | \textit{Self-employed} \hfill May 2019 -- May 2021
\begin{itemize}
  \item \href{https://www.github.com/karimn/takeup}{Built} a structural Bayesian model representing a behavior model, using data from a large-scale social experiment conducted in Kenya
  \item \href{https://github.com/karimn/covid-19-transmission}{Built} a hierarchical model to evaluate a scaled economic intervention program in Bangladesh
  \item Gaussian processes to model correlation in intervention effects by spatial proximity; hierarchical model of COVID-19 transmission heterogeneity under stay-at-home restrictions
  \item R \href{https://www.github.com/karimn/boundr}{package} for bounded identification of counterfactuals in DAGs under imperfect compliance
\end{itemize}

\medskip

\textbf{Economist} | \textit{Evidence Action} \hfill May 2014 -- May 2019
\begin{itemize}
  \item Designed and analyzed large-scale randomized evaluations for Evidence Action's flagship programs: deworming across East Africa and the at-scale replication of the No Lean Season seasonal-migration subsidy program in Bangladesh
  \item On No Lean Season, contributed to the analysis that failed to replicate the earlier pilot's effect on migration and consumption; the finding informed Evidence Action's decision to close the program --- an instance of rigorous evidence driving a hard organizational decision
  \item Work across projects spanned study design and power analysis, field data-collection supervision, and Bayesian and frequentist analysis for operational and policy audiences
\end{itemize}

\medskip

\textbf{Software Design Engineer in Test} | \textit{Microsoft}, Windows Debugging Tools \hfill Apr 1999 -- Mar 2007
\begin{itemize}
  \item Test framework and automation for Windows debugging tools; C/C++, systems-level testing
\end{itemize}

\section*{Projects}

\subsection*{Hierarchical state-space ``digital twin'' model for longitudinal tumor dynamics}

\textbf{Problem.} Oncology trials generate repeated tumor-size measurements that carry early predictive signal for survival, but standard analyses either reduce them to crude categorical response (RECIST) or fit patient-by-patient curves that ignore between-study structure. A Phase 3 initiation decision needs patient-level forecasts that borrow appropriately from hierarchically related trials.

\textbf{Method.} Built a Bayesian state-space ``digital twin'' model of log tumor size under a Stein-Fojo decay-plus-growth framework with a three-level hierarchy (population $\to$ trial $\to$ patient), non-centered parameterization, and Student-t priors on raw effects to absorb prior-data conflict in sparse strata (small subgroups). Each patient has a personalized posterior trajectory that borrows strength from their trial and the population. Implemented in Stan with modular include architecture; diagnostics included R-hat, ESS, E-BFMI, and prior-posterior conflict checks.

\textbf{Impact.} Extracts tumor-dynamics signal from Phase 2 and historical trial data to forecast Phase 3 outcomes at both population and patient levels, supporting Phase 3 initiation decisions on programs combining emerging trial evidence with real-world data.

\subsection*{Multistate illness-death survival model for competing-risk decomposition}

\textbf{Problem.} Progression-free survival bundles multiple distinct events --- disease progression and death --- under mechanisms that typically differ in hazard shape and covariate dependence. Conventional PFS analyses collapse this structure and can mislead trial-vs-historical comparisons when the event mix shifts between cohorts. And although much of progression is driven by underlying tumor-burden dynamics (SLD, PSA), important events are \textit{not} automatically explained by burden --- non-target lesion progression and clinical progression are clinician-adjudicated and not mechanically tied to the modelled burden trajectory.

\textbf{Method.} Built a Bayesian illness-death multistate model with three transitions (progression, death-without-pro\-gression, death-after-progression) \textbf{jointly fitted with the mechanistic tumor-dynamics model} (above). The patient-level latent tumor state enters the progression hazard directly, so burden-driven progression is captured mechanistically; the multistate layer adds the structure needed for progression events not driven by modelled burden and for competing death hazards. The same framework supports treating informative dropout --- non-administrative censoring and line-of-therapy switches --- as explicit transitions rather than ignorable censoring. Trial-level hierarchy; interval censoring for assessment-gap uncertainty; weighted likelihood for downstream real-world-data borrowing.

\textbf{Impact.} Makes competing-risk and informative-censoring structure explicit in survival comparisons, enabling defensible hazard-level comparisons between trial arms and historical or real-world cohorts where event mix and follow-up structure differ --- supporting Phase 3 initiation decisions on programs that combine emerging trial evidence with historical and real-world data.

\subsection*{Modular Bayesian inference infrastructure for clinical trial analyses}

\textbf{Problem.} Bayesian trial models accumulate components (tumor dynamics, multistate survival, covariate effects, data-borrowing schemes) that are repeatedly recombined for different trial analyses; hand-rolled Stan models drift in interface and diagnostics across analyses, making results hard to reproduce and compare.

\textbf{Method.} Designed a modular Stan architecture in which each sub-model (tumor regression, growth fraction, initial state, multistate hazards, measurement error, RWD borrowing) exposes a uniform 7-file interface --- flags, data, hyperparameters, transformed data, parameters, transformed parameters, priors. Wired into a \texttt{targets}-based R pipeline with standardized diagnostics reports, prior-sensitivity checks, and reproducible artifact storage.

\textbf{Impact.} Same components now feed multiple concurrent trial analyses with consistent diagnostics and priors; new model variants are composed rather than rewritten.

\subsection*{Leave-future-out cross-validation for hierarchical longitudinal--survival (``digital twin'') models}

\textbf{Problem.} Assessing the predictive accuracy of a patient-level tumor-dynamics and survival model requires scoring forecasts against data the model has not seen, but standard cross-validation conflates two kinds of held-out information --- \textit{future} observations from seen patients and \textit{unseen patients at seen times} --- and can make a model look better than it actually is at the forward-in-time forecasting task that matters for Phase 3 decision support. LFO-CV is well-established for time-series models (B\"{u}rkner, Gabry, Vehtari 2020) but is not routine for hierarchical longitudinal-plus-survival models.

\textbf{Method.} Adapted LFO-CV to a hierarchical state-space digital-twin model of tumor dynamics and survival: refit at successive truncation times and score out-of-sample predictive density on every held-out quantity --- longitudinal tumor-size visits \textit{and} survival/multistate events --- then summarize forecast performance as a function of prediction horizon. Primary workflow uses full-refit LFO-CV; PSIS-based approximate LFO is also supported for faster iteration. Built as a separate Stan variant sharing the base model's modular structure so validation tracks the production model exactly.

\textbf{Impact.} Produces horizon-stratified predictive-accuracy curves across both biomarker and survival outcomes --- a more honest answer to ``how well does the digital twin forecast?'' than aggregate LOO --- making forecast calibration directly auditable for Phase 3 decision reviews.

\subsection*{Agentic-workflow layer for Bayesian statistical modeling}

\textbf{Problem.} Bayesian trial-modeling work has two hard parts that off-the-shelf AI assistants don't address. First, the \textit{mathematical} layer: multistate survival with joint longitudinal submodels, non-centered hierarchical parameterizations, competing risks, and interval censoring produce dense derivations where a single person working by hand is both slow and error-prone. Second, the \textit{development-loop} layer: MCMC runs take hours to days, diagnostics and reparameterization require domain judgment, and state across the feedback loop is hard to carry.

\textbf{Method.} Designed and built an agentic-workflow layer on Claude Code, shared across the modeling team via an internal plugin marketplace. Core components: \textbf{collaborative mathematical derivation} --- working through the underlying statistical mathematics with Claude as a derivation partner before they become Stan code; a persistent memory system for carrying model-specific context across sessions; domain-specific sub-agents (Stan diagnostic triage, prior sensitivity review, \texttt{targets} pipeline explorer); MCP-integrated skills for pipeline operations and job management on Domino; and reusable skill and plugin templates that codify team conventions for Bayesian workflow. Version-controlled and iterated like any other team codebase.

\textbf{Impact.} Shifts the mathematical-modeling step from error-prone solo derivation to collaboratively-verified, fully-documented model specifications; compresses the hand-off between ``MCMC finished'' and ``I know what to do next,'' which is where most of the latency in Bayesian model development actually lives. Adopted across the modeling team for reproducible Stan development, MCMC diagnostics triage, and pipeline operations.

\section*{Publications}

\textbf{Peer-reviewed}

Barker, N., Davis, C.~A., L\'{o}pez-Pe\~{n}a, P., Mitchell, H., Mobarak, A.~M., Naguib, K., Reim\~{a}o, M., Shenoy, A., \& Vernot, C. ``Migration and Resilience during a Global Crisis.'' \textit{European Economic Review} (accepted).

\medskip

\textbf{Working papers}

Mitchell, H., Mobarak, A.~M., Naguib, K., Reim\~{a}o, M., \& Shenoy, A. ``External Validity and Implementation at Scale: Evidence from a Migration Loan Program in Bangladesh.''

\smallskip

Jee, E., Karing, A., \& Naguib, K. ``Optimal Incentives in the Presence of Social Norms: Experimental Evidence from Kenya.''

\medskip

\textbf{Preprint}

Naguib, K., Berch\'{e}, R., Li, L., Bevan, A., Khosla, S., Davies, J., \& Metcalfe, P. ``PIONEER: Bayesian Joint Modelling of Mechanistic Tumour Growth and Time-to-Event Endpoints for Dynamic Prediction of Ongoing Oncology Trials.'' \href{https://arxiv.org/abs/2607.17908}{arXiv:2607.17908}.

\end{document}
